\subsection{Full holomorphic endpoints of the fixed OS49 sphere families}
\begin{statement}{P3K.79}{Proved}
Retain an actual fixed-source OS49 sphere model of P3K.76--78,
with the given scalar, component Cartier line, finite cap of order
\(m=3k>0\), all fixed marked smooth order-one clutches and their full
root/line/annular data, and any periodic index-one lattice-basis fan \(P\).
Write \(D=n+m\sum_i\psi(\ell_i)=\pm1\) for its exact sphere datum.
Its full holomorphic automorphism group is the faithful physical
\(\mathbb C^*\) action. Every such automorphism descends to the identity
on the base and has identity linear part on the generic torus. The proof
uses the actual root and same-fan inversion maps to exclude the negative
linear part; no translation or source-descent assumption is imposed.

For a prescribed fixed-source scalar path of integer winding \(L\),
including the marked cap/line variations of P3K.44 and the fixed
smooth-site clutches of P3K.78, every specified holomorphic endpoint has
\(L=Nk\) for an integer \(N\). The explicit canonical endpoint and its
actual powers realize every such winding. All holomorphic endpoints for
that same path have oriented physical return
\[
                              N S_k(P)\in\mathbb Z/28,
\]
where \(S_k(P)\) is the actual all-fan return of P3K.77. The minimum
positive scalar winding in this endpoint scope is \(k\).
The maps retain the full finite group cochain, every line character,
and the original ordinary-clutch integers. This statement concerns the
fixed source and its stated data, with holomorphic endpoints; it does
not identify a different source-moduli path or a smooth endpoint that
has not been given a holomorphic map.
\end{statement}
\begin{proof}
\par

\medskip\noindent\textbf{FA.1. Recovering the base and fixing its three marked values}

\par\nobreak

Let \(\pi:X\to\mathbb P^1\) be the actual torus fibration. Its generic fibre T is a
compact complex two-torus, hence has a flat Kähler form \(\omega\). Since the
actual X is diffeomorphic to \(S^6\), \(H^2(X;\mathbb Z)=0\). If an effective divisor E
had a horizontal component, its restriction to a generic smooth fibre
would be a nonzero effective divisor \(E_T\). The line \(\mathcal O_X(E)\) has zero first
Chern class, so its restriction has zero first Chern class, whereas
\[
 \int_T c_1\bigl(\mathcal O_T(E_T)\bigr)\wedge\omega
                =\int_{E_T}\omega>0.                    \tag{FA.1}
\]
The equality is the divisor fundamental-class pairing; its last term is
strictly positive on each nonzero effective complex curve with its
multiplicity. This contradiction excludes every horizontal divisor.

The polar divisor of a meromorphic function on X is therefore vertical.
Its restriction to a generic compact connected T has no polar divisor;
removable singularities across the remaining codimension-two set make it
holomorphic, and hence constant. The local product over a regular base disc gives its meromorphic
descent there. At any special value choose a generic smooth point of a
reduced fibre component, outside the function's indeterminacy, and a
transverse analytic disc. In its local coordinate the base map is
\(\pi=u^e\) times a holomorphic unit. An extending e-th root of that unit
changes u so that \(\pi=u^e\) exactly. The pulled-back meromorphic function
has a convergent Laurent expansion with a finite principal part on
this disc. On the punctured disc it is constant on the generic fibres,
so this Laurent expansion is invariant under \(u\mapsto\zeta_e u\); every
nonzero exponent is divisible by e. Replacing \(u^e\) by the base
coordinate therefore gives the exact meromorphic extension across
that special value. This proves the descended function meromorphic
on all of \(\mathbb P^1\), and hence rational. Conversely every rational function
of pi pulls back. Thus \(\mathbb C(X)=\mathbb C(t)\), with this equality
implemented by \(\pi^*\). A biholomorphism F of X induces a Möbius transformation
h of \(\mathbb P^1\) and the exact equation \(\pi F=h\pi\). It holds first on regular fibres
and then everywhere by continuation.

The three nontrivial linear local monodromies are precisely \(A,B,C\) in
(OCS.3). Their characteristic polynomials are
\[
 \begin{aligned}
 \det(xI-A)&=(x-1)^2(x^2+x+1),\\
 \det(xI-B)&=(x-1)^2(x^2-x+1),\\
 \det(xI-C)&=(x-1)^4,\qquad C\ne I .
 \end{aligned}                                           \tag{FA.2}
\]
The actual finite root has its original effective order-three linear
monodromy A even when \(m=3k\); its ineffective subgroup is still retained in
the affine groupoid. Each added smooth order-one clutch has identity
linear monodromy. A holomorphic base map carries a positive local loop to
a positive local loop, so these conjugacy invariants make h fix 0,1 and
infinity separately. Hence h is the identity. This proves base preservation
for every biholomorphism without a prior source-preservation assumption.

\par

\medskip\noindent\textbf{FA.2. The complete integral linear centralizer}

\par\nobreak

A holomorphic map between compact complex tori lifts to their universal
covers. Its derivative is periodic and therefore constant, because its
entries descend to holomorphic functions on a compact torus. Every fibre
map is consequently affine. Its induced integral homology matrix is
locally constant over the regular base. Since F preserves the base, the
matrix X commutes with the actual three monodromies \(A,B,C\).

Commutation with B first gives the full block form
\[
 X=\begin{pmatrix}
 a&0&0&b\\0&p&q&0\\0&-q&p+q&0\\c&0&0&d
 \end{pmatrix}.
\]
Substitution in \(XA=AX\) gives \(b=q=0\) and \(a=p=d\). Consequently the complete
rational, and hence integral, centralizer is
\[
                  X=aI+c\,\delta\psi,
       \qquad \delta\psi=E_{41},\qquad \det X=a^4.       \tag{FA.3}
\]
An integral automorphism has \(a=\epsilon\in\{+1,-1\}\) and \(c\in\mathbb Z\). On a fixed
smooth torus its real linear map is complex linear. Subtracting
\(\epsilon I\) still gives a complex-linear map. If c were nonzero, the
result would have real rank one by (FA.3), whereas a nonzero complex-linear
map has even positive real rank. Thus \(c=0\). Every holomorphic automorphism
has linear part exactly \(+I\) or \(-I\). This computation uses the entire
monodromy module; it leaves no exceptional complex-multiplication case.

\par

\medskip\noindent\textbf{FA.3. Actual local inversion and the affine root/cusp square}

\par\nobreak

We use the same local origins and rigidified line coordinates as the
OCS groupoid. On every regular source chart, the degree-zero line has
constant factors of automorphy. In the corresponding additive-source,
multiplicative-line cover its group inverse is \((z,V)\mapsto(-z,V^{-1})\).
A source-translation line lift has the form \((z,V)\mapsto(z+s,cV)\): its ratio
with this flat lift is a holomorphic function on the compact elliptic
curve and hence is constant. These are actual translations of the same
commutative extension. Inversion conjugates the complete suspension B,
including its scalar, to \(B^{-1}\). This gives inversion of the local compact
torus quotient with its chosen origin.

Its extension at the retained II fibre has a direct Cartier formula.
Write \(i(x,y,t)=(x,-y,t)\), and let e be the canonical meromorphic frame of
\(M=\mathcal O(G-O)\), whose divisor is \(G-O\). On the ordinary source charts,
including the retained II disc and excluding the finite IV-star cap
and the selected cusp,
\[
 f=x-t^2,\qquad \operatorname{div}(f)=G+(-G)-2O,
 \qquad M\otimes i^*M\longrightarrow\mathcal O,
                 \quad e\otimes i^*e\longmapsto f.       \tag{FA.4}
\]
The nonzero-vector inverse is therefore literally
\[
                (p,V e(p))\longmapsto
                    \left(i(p),\frac{i^*e(p)}{f(p)V}\right).
                                                               \tag{FA.5}
\]
It is a map of the actual line bundles: the divisors in (FA.4) cancel the
apparent frame poles and zeros. Its square is the identity since \(f\circ i=f\).
At \(t=1\) the cusp has \(x=y=0\) and \(f=-1\), so the same map and its inverse extend
there. The rigidification arrow is explicit as well. In the source
parameter \(\tau=-x/y\) at O, the regular frame is \(\tau e\); since
\(\tau(i(p))=-\tau(p)\) and \(\tau^2 f\) tends to 1, the pairing (FA.4)
has value -1 on this regular frame at O. If the actual chosen origin
vector has coefficient \(v_i(q)\) in that frame, (FA.5) sends it to
\(-1/v_i(q)\). Compose (FA.5) with the extending physical scalar
\(-v_i(q)^2\). The resulting \(I_i\) fixes that exact origin and is its
group inverse. For the regular-frame origin \(v_i=1\) this factor is
literally -1. These are actual extending units, carried through the
line and all overlap maps; no parameter winding is discarded.

On the good root \(t=z^{3k},\ x=z^{4k}X,\ y=z^{6k}Y\), carry first the already proved corrected root-line map of F29/K37,
including its resolved-divisor correction and finite character. Its target
line is \(\mathcal O(G'-O')\), with \(G'=(z^{2k},1)\). The divisor pairing in (FA.4)
is, after removal of the displayed fibre power \(z^{4k}\), exactly
\[
                            f_{\rm root}=X-z^{2k}.        \tag{FA.6}
\]
The same regular-frame calculation at O supplies the identical
explicit factor \(-v_i(z)^2\) for its actual root origin. This defines
the origin-fixing inverse on the smooth root family. All original
finite line characters remain in the affine deck action. In particular
we do not claim that inversion commutes with an affine cap action: its
failure to do so is precisely the finite-group equation below. The
base-root lift of an actual base-preserving F is obtained by the literal
fibre product with \(z^m=t\) and its normalization; it fixes z. Its affine
linear part on the root torus is the same \(\epsilon I\).

At the selected cusp use the complete periods (OC.5), including every
coefficient function, and define on the toric covering
\[
                   I_c(Z,V,q)=(Z^{-1},V^{-1},q).           \tag{FA.7}
\]
It conjugates each full h and B period to its inverse. On the fan lattice
its map is \((v,h)\mapsto(-v,h)\). For the basis fan with triangles
\(\{0,u,v\}\) and \(\{u,v,u+v\}\), the negative of the first triangle, translated by
\(u+v\), is the second, and the negative of the second, translated by \(u+v\),
is the first. Thus this is an exact automorphism of the same infinite
periodic fan, compatible with every period and its full unit factors.
It extends on every toric chart, including non-semifree fans. This is an
actual local inverse map; no reflection to a different fan is substituted.
At every additional smooth clutch disc the ordinary local inverse above
is used, with the given local origin.

Suppose now the actual F had linear part \(-I\). Composing its local lifts
with these inversions gives identity-linear maps of the local models.
Their values at the chosen origin give extending translations \(s_i\). At
II and on the good root their line maps extend by (FA.4)--(FA.6) and the
actual Hom evaluation argument of OC.1. At the cusp, the origin lies in
the open central torus. Its image is reindexed by the literal h,B periods
using (OC.8); evaluation of the map and its inverse then gives the full
bounded units of (OC.7). Those units extend on every cone. These \(s_i\) are
therefore sections of the actual sheaf A in OCS, including the line and
all normal collars.

Write an ordinary overlap in precisely the OCS convention as
\[
                g_{ij}(x)=L_{ij}x+\theta_{ij},
                         \qquad F_i(x)=-x+s_i.
\]
The literal equation \(F_i g_{ij}=g_{ij}F_j\) is
\[
                       s_i-L_{ij}s_j=2\theta_{ij}.        \tag{FA.8}
\]
There is also the full finite-deck square. If the actual root arrow is
\(A_z(x)=L_D(z)x+a_D(z)\), from z to Dz, its square with F is
\[
                  s_{Dz}-L_D(z)s_z=2a_D(z).              \tag{FA.9}
\]
This retains the complete cyclic group of order m and every central line
character, including the ineffective subgroup. Taking all iterates of
(FA.9) retains the original m-th norm; it does not replace m by 3.
Equations (FA.8)--(FA.9) say exactly that twice the actual total translation
cocycle is the coboundary of these extending sections. Passing through
the same atlas exponential sequence (OCS.4), with its full integral
finite-group cochain, therefore gives
\[
                2(D,d\bmod k)=0
                  \quad\hbox{in }\mathbb Z\oplus\mathbb Z/k.     \tag{FA.10}
\]
But \(D=\pm1\) in every sphere datum under discussion. Its free component
\(2D\) is nonzero. Thus a linear \(-I\) automorphism cannot exist.

\par

\medskip\noindent\textbf{FA.4. The remaining group and endpoint consequence}

\par\nobreak

Every automorphism is now base preserving and identity linear, so OC.2--3
apply with their actual source/line and toric lifts. K78's explicit
clutch comparison transports that same group for the additional smooth
sites. The full calculation gives a rational lift \(f(t)U^j\), forces \(j=3a\),
and identifies it literally with \(cB_0^{-a}\); on the compact quotient it
is the physical scalar \(c\lambda^a\). Conversely every physical nonzero
scalar extends through every root and every toric chart, and acts
faithfully. Consequently
\[
                      \operatorname{Aut}_{\rm hol}(X)
                                =\mathbb C^*_{\rm physical}.     \tag{FA.11}
\]
A path from 1 to any c in this actual group is an isotopy of the full
complex model, its finite caps and its normal collars.

For a fixed closed family with two specified holomorphic endpoint maps,
the ratio of those endpoint maps is an automorphism of the same fixed
sphere member X. Equation (FA.11) supplies a literal isotopy between
them by physical scalar multiplication. Their inverse-endpoint physical
returns are therefore conjugate up to this isotopy and have the same
oriented clutched-sphere class. This conclusion concerns endpoints of the
same fixed data; it does not identify two different norm signs, different
source-moduli paths, or an arbitrary smooth endpoint with a holomorphic
one. No existence of an unspecified short endpoint is inferred.

The independent centralizer, characteristic-polynomial, Cartier, root,
fan and finite-cochain identities are recorded in
\path{code/part3/os49_full_automorphisms.py} and its captured output. The complete local inverse and root-groupoid argument have passed the independent audits recorded with this export.

\par

\medskip\noindent\textbf{FA.5. A prescribed scalar path with an arbitrary holomorphic endpoint}

\par\nobreak

Retain the actual fixed-source translation/line data of K44 and the fixed
smooth-site data of K78, with \(D=\pm1\). Let a cut scalar path have integer
winding L, and let an actual holomorphic endpoint between its two models
be specified. We now determine that endpoint's linear part without
assuming a translation marking.

On the ordinary source complement the endpoints have the same scalar
\(\lambda\) and hence the same literal quotient by \(B_\lambda\). Use this common
source/line-cover marking to compare their smooth compact tori. The
endpoint descends to the identity of the base by FA.1, and its integral
matrix commutes with the same \(A,B,C\). The rank-one argument of FA.2 again
leaves only \(+I\) and \(-I\). The analytic continuation of a chosen period
logarithm along the scalar path is explicitly
\[
                      S_L=I+L\delta\psi.                 \tag{FA.12}
\]
This preserves \(\psi\) and takes \(\ell_0\) to \(\ell_0+L\delta\). Thus every finite
norm vector and every transported ordinary vector retains its free
coordinate in (OCS.13). The actual total groupoid class at both endpoints
has free coordinate D. Its delta coordinate and every finite character
are retained, not equated by this assertion.

For an endpoint of linear part \(-I\), the literal affine square between
the two local atlases is now
\[
                    s_i-L_{ij}s_j
                         =\theta^{\rm end}_{ij}
                             +\theta^{\rm start}_{ij},    \tag{FA.13}
\]
with the identical sum in the full finite-deck square. The same local
extensions in FA.3 apply. Its integral free obstruction is therefore
\(D+D=2D\), nonzero. Thus the actual endpoint necessarily has linear part
\(+I\). The interval comparison of K78 consists of translations and hence
carries it to a translation endpoint of the canonical one-cap datum
\((m,D,d)\), with all scalar-origin and finite root terms retained.

K44, OC.15--19 now applies to this actual endpoint. Its forward root
square and the rational source-line divisor give
\[
                  3v-j=\frac{L D}{k},\qquad
                         L=N k\quad(N\in\mathbb Z).      \tag{FA.14}
\]
Here j is the actual source-translation coefficient and v its integral
cusp line order. The cap character, including its full norm, is the one
in OC.16--18; the source and line factors are not replaced by winding
numbers alone. The canonical OR65 endpoint and its powers realize every
\(L=Nk\), including negative N through the actual inverse map. K78 transports
those maps through the full family comparison.

For a given realized \(L=Nk\), compare any specified endpoint with that
transported canonical endpoint. Their ratio is an automorphism of the
same fixed sphere model; by (FA.11) it is an actual physical scalar.
Its isotopy removes that ratio from the inverse-endpoint physical
return. Thus every such holomorphic endpoint has the return
\[
                                N S_k(P).                \tag{FA.15}
\]
In particular the minimum positive scalar winding for this full
holomorphic endpoint scope is k. This conclusion is obtained from the
actual normalizer, affine and finite-root squares and the explicit
canonical endpoint; it is not a division of a class in Theta\_7. A family
with a changed source or an arbitrary smooth, nonholomorphic endpoint
has not been identified by this argument.

\end{proof}
